\documentclass[11pt]{article}
\usepackage{amsmath, amssymb, geometry, xcolor, enumitem}
\geometry{letterpaper, margin=1in}

\title{Homework 2.4 Exact Equations}
\author{Tashfeen Omran}
\date{September 2026}

\begin{document}

\maketitle
\tableofcontents
\newpage

\section{Part 1: Comprehensive Introduction, Context, and Prerequisites}

\subsection{Core Concepts}
The topic of this chapter is \textbf{Exact Differential Equations}. In multivariable calculus, if we have a function of two variables, $F(x, y)$, its total differential $dF$ is defined as:
\[ dF = \frac{\partial F}{\partial x} dx + \frac{\partial F}{\partial y} dy \]
A first-order differential equation of the form $M(x, y) dx + N(x, y) dy = 0$ is called an \textit{exact equation} if the expression on the left-hand side is the exact total differential of some underlying function $F(x, y)$. In other words, the equation is exact if there exists a function $F(x, y)$ such that:
\[ \frac{\partial F}{\partial x} = M(x, y) \quad \text{and} \quad \frac{\partial F}{\partial y} = N(x, y) \]
If the equation is exact, then the differential equation simplifies to $dF = 0$. Integrating both sides yields the general implicit solution:
\[ F(x, y) = C \]

\subsection{Intuition and Derivation}
How do we know if a given equation is exact without guessing the function $F(x, y)$? We rely on \textbf{Clairaut's Theorem} from multivariable calculus, which states that for well-behaved functions, mixed second partial derivatives are equal. That is:
\[ \frac{\partial^2 F}{\partial y \partial x} = \frac{\partial^2 F}{\partial x \partial y} \]
If we assume $M = \frac{\partial F}{\partial x}$ and $N = \frac{\partial F}{\partial y}$, then taking the derivative of $M$ with respect to $y$ and $N$ with respect to $x$ gives:
\[ \frac{\partial M}{\partial y} = \frac{\partial^2 F}{\partial y \partial x} \quad \text{and} \quad \frac{\partial N}{\partial x} = \frac{\partial^2 F}{\partial x \partial y} \]
By Clairaut's Theorem, if the equation is exact, these two partial derivatives must be strictly equal. This gives us our fundamental test for exactness:
\[ \frac{\partial M}{\partial y} = \frac{\partial N}{\partial x} \]

\subsection{Historical Context and Motivation}
The concept of exact differentials was rigorously developed in the 18th century by mathematicians like Leonhard Euler and Alexis Clairaut. They were motivated by problems in classical mechanics and thermodynamics, specifically the concept of "conservative vector fields." Physicists needed to know if the work done by a force field was independent of the path taken by an object. If a differential expression was exact, it meant it possessed a "potential function" (analogous to our $F(x,y)$), meaning energy was conserved, and the system was vastly easier to solve.

\subsection{Key Formulas}
\begin{itemize}
    \item \textbf{Standard Form:} $M(x, y) dx + N(x, y) dy = 0$
    \item \textbf{Test for Exactness:} $M_y = N_x$ (where $M_y = \frac{\partial M}{\partial y}$ and $N_x = \frac{\partial N}{\partial x}$)
    \item \textbf{Recovering $F(x,y)$:} $F(x,y) = \int M(x,y) dx + g(y)$ or $F(x,y) = \int N(x,y) dy + h(x)$
    \item \textbf{General Solution:} $F(x, y) = C$
\end{itemize}

\subsection{Prerequisites}
To master this topic, you must be proficient in:
\begin{itemize}
    \item \textbf{Partial Differentiation:} Taking a derivative with respect to one variable while treating all other variables as constants.
    \item \textbf{Partial Integration:} Integrating an expression with respect to $x$ while treating $y$ as a constant. The "constant of integration" will be a function of the other variable, $g(y)$.
    \item \textbf{Implicit Functions:} Understanding how to express solutions implicitly as $F(x,y) = C$ rather than explicitly as $y = f(x)$.
\end{itemize}

\newpage
\section{Part 2: Detailed Homework Solutions}

\subsection*{Problem 1}
\textbf{Question:} Classify the equation as separable, linear, exact, or none of these. Note that it is possible for the equation to have more than one classification.
\[ \left(x^{\frac{9}{4}} - 8y\right) dx + x dy = 0 \]
\textbf{Solution:} Let us check all three classifications.
1. \textbf{Exactness:} 
Here, $M = x^{9/4} - 8y$ and $N = x$.
\[ M_y = -8, \quad N_x = 1 \]
Since $M_y \neq N_x$, the equation is \textbf{not exact}.

2. \textbf{Linearity:}
Let us rearrange the equation into the standard linear form, $\frac{dy}{dx} + P(x)y = Q(x)$.
Divide the entire equation by $x dx$:
\[ \frac{x^{9/4}}{x} - \frac{8y}{x} + \frac{dy}{dx} = 0 \]
\[ \frac{dy}{dx} - \frac{8}{x}y = -x^{5/4} \]
This perfectly matches the linear form with $P(x) = -\frac{8}{x}$ and $Q(x) = -x^{5/4}$. Thus, the equation is \textbf{linear}.

3. \textbf{Separability:}
An equation is separable if it can be written as $f(y) dy = g(x) dx$. Looking at our linear form, we have a $y$ term mixed via subtraction with an $x$ term, and we cannot factor this into a product of purely $x$ and purely $y$ functions. Thus, it is \textbf{not separable}.

\textbf{Final Answer:} Select \textbf{B. Linear}.

\vspace{0.5cm}
\subsection*{Problem 2}
\textbf{Question:} Classify the equation as separable, linear, exact, or none of these. 
\[ \sqrt{2y - y^2} dx + \left(7 + 4x - x^2\right) dy = 0 \]
\textbf{Solution:} Let us analyze the classifications.
1. \textbf{Exactness:}
$M = \sqrt{2y - y^2}$ and $N = 7 + 4x - x^2$.
\[ M_y = \frac{1}{2\sqrt{2y-y^2}}(2 - 2y) = \frac{1-y}{\sqrt{2y-y^2}}, \quad N_x = 4 - 2x \]
Since $M_y \neq N_x$, it is \textbf{not exact}.

2. \textbf{Linearity:}
The presence of $\sqrt{y}$ immediately disqualifies the equation from being linear, as linear equations require $y$ and its derivatives to be to the first power and not trapped inside non-linear functions.

3. \textbf{Separability:}
Let's try to isolate $x$ terms with $dx$ and $y$ terms with $dy$. 
\[ \sqrt{2y - y^2} dx = -(7 + 4x - x^2) dy \]
Dividing both sides by $\sqrt{2y - y^2}$ and $(7 + 4x - x^2)$:
\[ \frac{dx}{7 + 4x - x^2} = \frac{-dy}{\sqrt{2y - y^2}} \]
We have successfully separated the variables! Thus, the equation is \textbf{separable}.

\textbf{Final Answer:} Select \textbf{A. Separable}.

\vspace{0.5cm}
\subsection*{Problem 3}
\textbf{Question:} Classify the equation as separable, linear, exact, or none of these.
\[ \left(5x e^{5xy} + 9y\right) dx + \left(5y e^{5xy} + 9x\right) dy = 0 \]
\textbf{Solution:} 
1. \textbf{Exactness:}
$M = 5x e^{5xy} + 9y$ and $N = 5y e^{5xy} + 9x$.
Applying the product rule and chain rule to $M_y$:
\[ M_y = 5x \left(e^{5xy} \cdot 5x\right) + 9 = 25x^2 e^{5xy} + 9 \]
Applying the rules to $N_x$:
\[ N_x = 5y \left(e^{5xy} \cdot 5y\right) + 9 = 25y^2 e^{5xy} + 9 \]
Since $25x^2 e^{5xy} + 9 \neq 25y^2 e^{5xy} + 9$, it is \textbf{not exact}. 
\textit{(Note: If the coefficients were swapped such that $M$ contained $5ye^{5xy}$ and $N$ contained $5xe^{5xy}$, it would have been exact.)}

2. \textbf{Linearity:}
The term $e^{5xy}$ involves the dependent variable $y$ in the exponent, which is highly non-linear. Thus, it is \textbf{not linear}.

3. \textbf{Separability:}
The exponential term $e^{5xy}$ inextricably binds $x$ and $y$ together through addition and multiplication with other terms. It is impossible to factor this equation into the form $f(x) dx = g(y) dy$. Thus, it is \textbf{not separable}.

\textbf{Final Answer:} Select \textbf{D. None of these}.

\vspace{0.5cm}
\subsection*{Problem 4}
\textbf{Question:} Determine whether the equation is exact. If it is, then solve it.
\[ \left(3x^2 y + 7\right) dx + \left(x^3 - 4\right) dy = 0 \]
\textbf{Solution:} 
\textbf{Step 1: Test for exactness.}
\[ M(x,y) = 3x^2 y + 7 \implies M_y = 3x^2 \]
\[ N(x,y) = x^3 - 4 \implies N_x = 3x^2 \]
Since $M_y = N_x$, the equation is \textbf{exact}.

\textbf{Step 2: Integrate $M$ with respect to $x$.}
\[ F(x, y) = \int M(x, y) dx = \int \left(3x^2 y + 7\right) dx = x^3 y + 7x + g(y) \]

\textbf{Step 3: Differentiate $F(x,y)$ with respect to $y$ and set equal to $N$.}
\[ F_y = x^3 + g'(y) \]
We know $F_y = N(x,y)$, so:
\[ x^3 + g'(y) = x^3 - 4 \]
\[ g'(y) = -4 \]

\textbf{Step 4: Integrate to find $g(y)$.}
\[ g(y) = \int -4 \, dy = -4y \]
\textit{(We do not need to add a constant here as it gets absorbed into the final constant $C$).}

\textbf{Step 5: Write the final implicit solution.}
Substitute $g(y)$ back into $F(x,y)$:
\[ F(x, y) = x^3 y + 7x - 4y \]
\textbf{Final Answer:} Select \textbf{A}. The equation is exact, and the implicit solution is $\mathbf{x^3 y + 7x - 4y = C}$.

\vspace{0.5cm}
\subsection*{Problem 5}
\textbf{Question:} Determine whether the equation is exact. If it is, then solve it.
\[ e^t (7y - 5t) dt + \left(6 + 7e^t\right) dy = 0 \]
\textbf{Solution:} 
\textbf{Step 1: Test for exactness.} (Here, our variables are $t$ and $y$).
Let's expand $M(t,y)$:
\[ M(t,y) = 7y e^t - 5t e^t \implies M_y = 7e^t \]
\[ N(t,y) = 6 + 7e^t \implies N_t = 7e^t \]
Since $M_y = N_t$, the equation is \textbf{exact}.

\textbf{Step 2: Integrate $N$ with respect to $y$ (it's simpler than integrating $M$).}
\[ F(t,y) = \int N(t,y) dy = \int \left(6 + 7e^t\right) dy = 6y + 7y e^t + h(t) \]

\textbf{Step 3: Differentiate $F(t,y)$ with respect to $t$ and set equal to $M$.}
\[ F_t = 7y e^t + h'(t) \]
Set this equal to $M(t,y)$:
\[ 7y e^t + h'(t) = 7y e^t - 5t e^t \]
\[ h'(t) = -5t e^t \]

\textbf{Step 4: Integrate to find $h(t)$ using Integration by Parts.}
We must evaluate $\int -5t e^t dt$. Let $u = -5t$ and $dv = e^t dt$.
Then $du = -5 dt$ and $v = e^t$.
\[ \int u \, dv = uv - \int v \, du = -5t e^t - \int e^t (-5) dt = -5t e^t + 5e^t \]
Thus, $h(t) = -5t e^t + 5e^t$.

\textbf{Step 5: Write the final implicit solution.}
Substitute $h(t)$ into $F(t,y)$:
\[ F(t, y) = 6y + 7y e^t - 5t e^t + 5e^t \]
\textbf{Final Answer:} Select \textbf{A}. The equation is exact, and the implicit solution is $\mathbf{6y + 7ye^t - 5te^t + 5e^t = C}$.

\vspace{0.5cm}
\subsection*{Problem 6}
\textbf{Question:} Determine whether the equation is exact. If it is, then solve it.
\[ \left( \frac{3}{y} \right) dx - \left( 6y - \frac{3x}{y^2} \right) dy = 0 \]
\textbf{Solution:} 
\textbf{Step 1: Distribute the negative sign to define $N(x,y)$ properly.}
The standard form is $M dx + N dy = 0$. Therefore:
\[ M(x, y) = \frac{3}{y} \]
\[ N(x, y) = -\left( 6y - \frac{3x}{y^2} \right) = -6y + \frac{3x}{y^2} \]

\textbf{Step 2: Calculate partial derivatives to test for exactness.}
\[ \frac{\partial M}{\partial y} = -\frac{3}{y^2} \]
\[ \frac{\partial N}{\partial x} = \frac{\partial}{\partial x}\left( -6y + \frac{3x}{y^2} \right) = \frac{3}{y^2} \]
Since $-\frac{3}{y^2} \neq \frac{3}{y^2}$, we conclude that $M_y \neq N_x$. The equation fails the exactness test.
\textit{(Note: If the original equation had a `+` sign between the terms instead of a `-` sign, it would have been exact. However, based strictly on the equation as written, it is not.)}

\textbf{Final Answer:} Select \textbf{B. The equation is not exact.}

\vspace{0.5cm}
\subsection*{Problem 7}
\textbf{Question:} Solve the initial value problem. Type an exact answer in terms of $\pi$.
\[ \left( \frac{1}{x} + 8y^2 x \right) dx + \left(8y x^2 - \cos y\right) dy = 0, \quad y(1) = \pi \]
\textbf{Solution:} 
\textbf{Step 1: Test for exactness.}
\[ M_y = 16xy \]
\[ N_x = 16xy \]
Exact!

\textbf{Step 2: Find the potential function $F(x,y)$.}
\[ F(x, y) = \int M \, dx = \int \left( \frac{1}{x} + 8y^2 x \right) dx = \ln|x| + 4y^2 x^2 + g(y) \]
\textit{(Because our initial condition has $x = 1 > 0$, we can safely drop the absolute value bars and write $\ln x$).}

\textbf{Step 3: Find $g(y)$.}
\[ F_y = 8y x^2 + g'(y) \]
Set $F_y = N$:
\[ 8y x^2 + g'(y) = 8y x^2 - \cos y \implies g'(y) = -\cos y \]
\[ g(y) = \int -\cos y \, dy = -\sin y \]
Thus, the implicit general solution is $\ln x + 4x^2 y^2 - \sin y = C$.

\textbf{Step 4: Apply the initial condition $y(1) = \pi$.}
Plug in $x = 1$ and $y = \pi$:
\[ \ln(1) + 4(1)^2 (\pi)^2 - \sin(\pi) = C \]
Since $\ln(1) = 0$ and $\sin(\pi) = 0$, we have:
\[ 0 + 4\pi^2 - 0 = C \implies C = 4\pi^2 \]

\textbf{Final Answer:} $\mathbf{\ln x + 4x^2 y^2 - \sin y = 4\pi^2}$

\vspace{0.5cm}
\subsection*{Problem 8}
\textbf{Question:} Solve the initial value problem.
\[ \left( 7y e^{7xy} - \frac{9}{y} \right) dx + \left( 7x e^{7xy} + \frac{9x}{y^2} \right) dy = 0, \quad y(0) = 2 \]
\textbf{Solution:}
\textbf{Step 1: Test for exactness.}
\[ M_y = \frac{\partial}{\partial y}\left(7y e^{7xy} - 9y^{-1}\right) = 7\left(1 \cdot e^{7xy} + y \cdot e^{7xy} \cdot 7x\right) + 9y^{-2} = 7e^{7xy} + 49xy e^{7xy} + \frac{9}{y^2} \]
\[ N_x = \frac{\partial}{\partial x}\left(7x e^{7xy} + 9xy^{-2}\right) = 7\left(1 \cdot e^{7xy} + x \cdot e^{7xy} \cdot 7y\right) + \frac{9}{y^2} = 7e^{7xy} + 49xy e^{7xy} + \frac{9}{y^2} \]
$M_y = N_x$, so it is exact.

\textbf{Step 2: Find the potential function $F(x,y)$.}
\[ F(x, y) = \int M \, dx = \int \left( 7y e^{7xy} - \frac{9}{y} \right) dx \]
To integrate $7y e^{7xy}$ with respect to $x$, treat $y$ as a constant. The integral is $e^{7xy}$ because the derivative of $e^{7xy}$ wrt $x$ is $e^{7xy} \cdot 7y$.
\[ F(x, y) = e^{7xy} - \frac{9x}{y} + g(y) \]

\textbf{Step 3: Find $g(y)$.}
\[ F_y = 7x e^{7xy} - 9x(-1)y^{-2} + g'(y) = 7x e^{7xy} + \frac{9x}{y^2} + g'(y) \]
Set $F_y = N$:
\[ 7x e^{7xy} + \frac{9x}{y^2} + g'(y) = 7x e^{7xy} + \frac{9x}{y^2} \implies g'(y) = 0 \implies g(y) = 0 \]
General solution: $e^{7xy} - \frac{9x}{y} = C$.

\textbf{Step 4: Apply the initial condition $y(0) = 2$.}
\[ e^{7(0)(2)} - \frac{9(0)}{2} = C \implies e^0 - 0 = C \implies C = 1 \]

\textbf{Final Answer:} $\mathbf{e^{7xy} - \frac{9x}{y} = 1}$

\vspace{0.5cm}
\subsection*{Problem 9}
\textbf{Question:} Determine whether the equation is exact. If it is, then solve it.
\[ \left[ 2y \cos(xy) + \frac{3}{\sqrt{1-x^2}} \right] dx + \left[ 2x \cos(xy) - 5y^{-1/5} \right] dy = 0 \]
\textbf{Solution:}
\textbf{Step 1: Test for exactness.}
\[ M_y = 2\cos(xy) + 2y(-\sin(xy) \cdot x) + 0 = 2\cos(xy) - 2xy \sin(xy) \]
\[ N_x = 2\cos(xy) + 2x(-\sin(xy) \cdot y) - 0 = 2\cos(xy) - 2xy \sin(xy) \]
Exact!

\textbf{Step 2: Find $F(x,y)$.}
\[ F(x,y) = \int M \, dx = \int \left( 2y \cos(xy) + \frac{3}{\sqrt{1-x^2}} \right) dx = 2\sin(xy) + 3\arcsin(x) + g(y) \]

\textbf{Step 3: Find $g(y)$.}
\[ F_y = 2x \cos(xy) + g'(y) \]
Set $F_y = N$:
\[ 2x \cos(xy) + g'(y) = 2x \cos(xy) - 5y^{-1/5} \implies g'(y) = -5y^{-1/5} \]
\[ g(y) = \int -5y^{-1/5} \, dy = -5 \left( \frac{y^{4/5}}{4/5} \right) = -5 \cdot \frac{5}{4} y^{4/5} = -\frac{25}{4}y^{4/5} \]

\textbf{Step 4: Final solution.}
Substitute $g(y)$ to get $F(x,y) = C$.
\textbf{Final Answer:} Select \textbf{A}. The equation is exact, and the implicit solution is $\mathbf{2\sin(xy) + 3\arcsin(x) - \frac{25}{4}y^{4/5} = C}$.

\vspace{0.5cm}
\subsection*{Problem 10}
\textbf{Question:} Solve the initial value problem.
\[ (\tan(9y) - 1) dx + \left( 9x \sec^2(9y) + \frac{1}{y} \right) dy = 0, \quad y(0) = 1 \]
\textbf{Solution:}
\textbf{Step 1: Test for exactness.}
\[ M_y = \frac{\partial}{\partial y}(\tan(9y) - 1) = 9\sec^2(9y) \]
\[ N_x = \frac{\partial}{\partial x}\left(9x \sec^2(9y) + \frac{1}{y}\right) = 9\sec^2(9y) \]
Exact!

\textbf{Step 2: Find $F(x,y)$.}
\[ F(x,y) = \int (\tan(9y) - 1) dx = x\tan(9y) - x + g(y) \]

\textbf{Step 3: Find $g(y)$.}
\[ F_y = x(9\sec^2(9y)) + g'(y) = 9x \sec^2(9y) + g'(y) \]
Set $F_y = N$:
\[ 9x \sec^2(9y) + g'(y) = 9x \sec^2(9y) + \frac{1}{y} \implies g'(y) = \frac{1}{y} \implies g(y) = \ln|y| \]
Since our initial condition dictates $y(0) = 1$, we are in a domain where $y > 0$, so we can drop the absolute value bars: $g(y) = \ln y$.
General solution: $x\tan(9y) - x + \ln y = C$.

\textbf{Step 4: Apply IVP condition.}
Plug in $x = 0$ and $y = 1$:
\[ (0)\tan(9) - 0 + \ln(1) = C \implies 0 - 0 + 0 = C \implies C = 0 \]

\textbf{Final Answer:} $\mathbf{x\tan(9y) - x + \ln y = 0}$

\newpage
\section{Part 3: In-Depth Analysis of Problems and Techniques}

\subsection{A) Problem Types and General Approach}
We can categorize the homework into three main problem structures:

\begin{enumerate}
    \item \textbf{Classification Problems (Problems 1, 2, 3, 6):}
    \begin{itemize}
        \item \textbf{Approach:} For these problems, solving isn't the primary goal. You must run the equation through a series of "tests." 
        \item \textit{Exact test:} Calculate $M_y$ and $N_x$.
        \item \textit{Linear test:} Rearrange algebraically to $\frac{dy}{dx} + P(x)y = Q(x)$. Ensure $y$ isn't trapped in a root or exponential.
        \item \textit{Separable test:} Attempt to cross-multiply and factor the equation so that all $x$'s and $dx$ are on one side, and all $y$'s and $dy$ are on the other.
    \end{itemize}
    
    \item \textbf{Standard Exact Differential Equations (Problems 4, 5, 9):}
    \begin{itemize}
        \item \textbf{Approach:} The primary goal is to find the potential function $F(x,y)$. You integrate $M$ with respect to $x$ to get an expression with an unknown $g(y)$. You then differentiate that expression with respect to $y$, set it equal to $N$, and solve for $g(y)$. You assemble the final pieces into the form $F(x,y) = C$.
    \end{itemize}

    \item \textbf{Initial Value Problems (IVPs) (Problems 7, 8, 10):}
    \begin{itemize}
        \item \textbf{Approach:} These are solved identically to the Standard Exact equations, but they have an extra arithmetic step at the end. Once you have the implicit general solution $F(x,y) = C$, you plug in the provided initial conditions (an $(x,y)$ coordinate) to solve for the specific numerical value of $C$.
    \end{itemize}
\end{enumerate}

\subsection{B) Key Algebraic and Calculus Manipulations}
\begin{itemize}
    \item \textbf{Partial Differentiation with Product/Chain Rule (Problem 3, 8, 9):} Differentiating multivariable functions requires strict adherence to treating the "other" variable as a constant. For example, in Problem 9, $\frac{\partial}{\partial y}[2y \cos(xy)]$ requires the product rule where $u = 2y$ and $v = \cos(xy)$. This resulted in $2\cos(xy) - 2xy \sin(xy)$.
    \item \textbf{Integration by Parts (Problem 5):} To integrate $g'(t) = -5t e^t$, standard substitution fails. We applied Integration by Parts ($\int u dv = uv - \int v du$) by choosing $u = -5t$ to reduce its degree through differentiation, making the remaining integral trivial.
    \item \textbf{Fractional Exponent Power Rule (Problem 9):} We utilized the standard power rule for integration $\int x^n dx = \frac{x^{n+1}}{n+1}$ on the tricky fraction exponent: $\int -5y^{-1/5} dy = -5 \frac{y^{4/5}}{4/5}$. Multiplying by the reciprocal $\frac{5}{4}$ streamlined the algebra into $-\frac{25}{4}y^{4/5}$.
    \item \textbf{Dropping Absolute Value Bars in Logs (Problems 7, 10):} Integrating $\frac{1}{y}$ results strictly in $\ln|y|$. However, we observed the given Initial Value Condition (e.g., $y(0) = 1$). Because the initial condition establishes that the function starts in the positive domain, and logarithmic functions are continuous over their domain, we can safely drop the absolute value bars and write $\ln(y)$.
\end{itemize}

\newpage
\section{Part 4: "Cheatsheet" and Tips for Success}

\subsection{Summary of Most Important Formulas}
\begin{itemize}
    \item \textbf{Test for Exactness:} $\frac{\partial M}{\partial y} = \frac{\partial N}{\partial x}$
    \item \textbf{Potential Function Recovery:} $F(x,y) = \int M dx + g(y) = C$ 
    \item \textbf{Finding the missing $g(y)$:} $g'(y) = N(x,y) - \frac{\partial}{\partial y} \int M dx$
\end{itemize}

\subsection{Tricks and Shortcuts}
\begin{itemize}
    \item \textbf{Choose the simpler integral:} If $M$ is extremely complicated but $N$ is simple, do not integrate $M$ with respect to $x$ first! You can just as easily start by integrating $N$ with respect to $y$, which gives $F(x,y) = \int N dy + h(x)$. Then differentiate with respect to $x$ and set equal to $M$ to find $h'(x)$. (We used this strategy intuitively in Problem 5!).
    \item \textbf{Visual Exactness Check:} Sometimes, you can spot exactness by recognizing the reverse product rule. For example, $y dx + x dy$ is immediately exact because it is the exact differential $d(xy)$.
\end{itemize}

\subsection{Common Pitfalls, Mistakes, and Traps}
\begin{itemize}
    \item \textbf{The Missing Minus Sign Trap:} The standard form is $M dx + N dy = 0$. If an equation is written with a minus sign, such as $M dx - (\text{expression}) dy = 0$, you \textbf{must} distribute that minus sign into the expression before defining $N$. Failing to do this guarantees you will calculate exactness incorrectly (as analyzed in Problem 6).
    \item \textbf{Forgetting the $+ C$ in the General Solution:} Exact equations do not yield explicit functions $y = f(x)$. They yield implicit relations $F(x,y) = C$. If you just write down the expression $F(x,y)$ without setting it equal to $C$, your answer is not an equation and is conceptually meaningless.
    \item \textbf{Mixing up variables during partials:} It is incredibly common to accidentally take $M_x$ and $N_y$. Remember the "criss-cross" rule: differentiate the term attached to $dx$ with respect to $y$, and the term attached to $dy$ with respect to $x$.
\end{itemize}

\newpage
\section{Part 5: Conceptual Synthesis and The "Big Picture"}

\subsection{A) Thematic Connections}
The core theme of exact differential equations is \textbf{Path Independence and Potential Functions}. This ties directly back to Multivariable Calculus (Calculus III). 
When you check if $M_y = N_x$, you are effectively checking if the 2D vector field $\mathbf{F} = \langle M, N \rangle$ has a "curl" of zero. If it does, the vector field is conservative. 
Finding the solution $F(x,y) = C$ is exactly the same mathematical process as finding the scalar potential function $f(x,y)$ such that $\nabla f = \mathbf{F}$. The lines where $F(x,y) = C$ represent the equipotential contour lines of that scalar field.

\subsection{B) Forward and Backward Links}
\begin{itemize}
    \item \textbf{Backward Links:} This topic relies entirely on the Chain Rule and Implicit Differentiation from Calculus I, as well as Partial Differentiation from Calculus III. We rely on Clairaut's theorem to prove that if the mixed partials are equal, an underlying function exists.
    \item \textbf{Forward Links:} Mastery of exact differentials is mandatory for advanced physics and thermodynamics (e.g., Maxwell's Relations). Furthermore, it lays the groundwork for evaluating Line Integrals using the Fundamental Theorem of Line Integrals. If a differential equation is exact, integrating it over a curve depends only on the endpoints, not the path taken!
\end{itemize}

\newpage
\section{Part 6: Real-World Application and Modeling}

\subsection{A) Concrete Scenarios in Economics, Finance, and Stats}
\begin{enumerate}
    \item \textbf{Economics (Indifference Curves):} In microeconomics, a consumer's total utility $U(x,y)$ derived from consuming goods $X$ and $Y$ produces marginal utilities $MU_x$ and $MU_y$. The differential equation $MU_x dx + MU_y dy = 0$ represents an indifference curve (a curve where total utility is constant). Testing for exactness checks if a given set of marginal utilities mathematically stems from a valid, consistent total utility function.
    \item \textbf{Quantitative Finance (Stochastic Calculus):} In financial modeling, stock prices often follow Itô processes. When deriving pricing PDEs (like the Vasicek model for interest rates), quants often perform variable substitutions that reduce complex PDEs into solvable ordinary differential equations. Recognizing these ODEs as exact (or using integrating factors to force them to be exact) is how we arrive at closed-form pricing formulas for European derivatives.
    \item \textbf{Probability Theory (Reconstructing Joint CDFs):} In multivariate statistics, if a researcher has derived the partial marginal change rates of a joint Cumulative Distribution Function (CDF), they can set up an exact equation. Solving the exact equation integrates the partial information back into the complete joint CDF, mapping the holistic probability space.
\end{enumerate}

\subsection{B) Model Problem Setup: The Consumer Utility Modeler}
\textbf{Scenario:} An economist has empirically measured how a consumer's satisfaction changes based on tiny increments of purchasing apples ($x$) and bananas ($y$). 
The marginal utility of an apple is modeled as $M(x,y) = 3x^2y + 2y^2$. 
The marginal utility of a banana is modeled as $N(x,y) = x^3 + 4xy$.
The economist needs to find the overarching Total Utility function $U(x,y)$ to map out the consumer's indifference curves across the market.

\textbf{Mathematical Setup:}
\begin{itemize}
    \item \textbf{Variables:} $x$ = Quantity of apples, $y$ = Quantity of bananas.
    \item \textbf{Given:} $M(x,y) = 3x^2y + 2y^2$ and $N(x,y) = x^3 + 4xy$.
    \item \textbf{Differential Equation:} The indifference curve implies no change in utility, so $dU = 0$.
    \[ (3x^2y + 2y^2) dx + (x^3 + 4xy) dy = 0 \]
    \item \textbf{The Task to Solve:} The economist must first verify the model is structurally sound by checking $M_y = N_x$. If $3x^2 + 4y = 3x^2 + 4y$ (which it does), the equation is exact. The economist would then integrate $M$ with respect to $x$ and solve for $g(y)$ to yield the total utility function $U(x,y) = C$, which represents the full map of consumer preferences!
\end{itemize}

\newpage
\section{Part 7: Common Variations and Untested Concepts}

The most significant concept taught in standard exact equations curricula that was completely missing from your homework is \textbf{Integrating Factors for Non-Exact Equations}.

\textbf{The Concept:} 
Often, you will find an equation where $M_y \neq N_x$. While your homework taught you to just label these "Not Exact" and move on, in reality, we can often \textit{force} them to become exact by multiplying the entire equation by a special function called an Integrating Factor, denoted as $\mu$.

\textbf{How it works:}
\begin{itemize}
    \item If $\frac{M_y - N_x}{N}$ evaluates to a function of \textbf{only} $x$, say $P(x)$, then the integrating factor is $\mu(x) = e^{\int P(x) dx}$.
    \item If $\frac{N_x - M_y}{M}$ evaluates to a function of \textbf{only} $y$, say $Q(y)$, then the integrating factor is $\mu(y) = e^{\int Q(y) dy}$.
\end{itemize}
Multiplying the original non-exact equation by $\mu$ will yield a new equation that perfectly passes the exactness test!

\textbf{Worked Example:}
Solve $\left(2y^2 + 3x\right) dx + 2xy dy = 0$.

\textbf{Step 1:} Test exactness. $M_y = 4y$ and $N_x = 2y$. $4y \neq 2y$, so it is not exact.

\textbf{Step 2:} Find an Integrating Factor.
Let's check the $x$-only test:
\[ \frac{M_y - N_x}{N} = \frac{4y - 2y}{2xy} = \frac{2y}{2xy} = \frac{1}{x} \]
Since this depends only on $x$, our integrating factor is:
\[ \mu(x) = e^{\int \frac{1}{x} dx} = e^{\ln x} = x \]

\textbf{Step 3:} Multiply original equation by $x$.
\[ \left(2xy^2 + 3x^2\right) dx + 2x^2 y dy = 0 \]
Now, check exactness again. New $M_y = 4xy$. New $N_x = 4xy$. It is now exact and solvable using standard methods!

\newpage
\section{Part 8: Advanced Diagnostic Testing: "Find the Flaw"}

Below are 5 problems with subtle, common mistakes. Your task as a grader is to find the flaw, explain it, and provide the correct step.

\vspace{0.5cm}
\noindent \textbf{Problem 1: The Double Counting Error}
\newline \textbf{Equation:} $(2xy) dx + (x^2 + 3y^2) dy = 0$
\newline \textbf{Flawed Solution:} 
$M = 2xy \implies M_y = 2x$. $N = x^2 + 3y^2 \implies N_x = 2x$. Exact.
Integrate $M$ wrt $x$: $\int 2xy dx = x^2 y$.
Integrate $N$ wrt $y$: $\int (x^2 + 3y^2) dy = x^2 y + y^3$.
Combine all terms to form the final equation:
$F(x,y) = x^2 y + x^2 y + y^3 = 2x^2 y + y^3 = C$.
\newline \textbf{Find the Flaw:}
\vspace{1cm} % Space for user to think

\noindent \textbf{Problem 2: The Variable Mix-up}
\newline \textbf{Equation:} $(3x^2 + 2y) dx + (2x + 3y^2) dy = 0$
\newline \textbf{Flawed Solution:} 
Let $M = 3x^2 + 2y$ and $N = 2x + 3y^2$. 
To test for exactness, we take the derivatives:
$M_x = 6x$ and $N_y = 6y$.
Since $6x \neq 6y$, the equation is NOT exact and cannot be solved with this method.
\newline \textbf{Find the Flaw:}
\vspace{1cm}

\noindent \textbf{Problem 3: The Integration Amnesia}
\newline \textbf{Equation:} $(2xy + 1) dx + (x^2 + \cos y) dy = 0$
\newline \textbf{Flawed Solution:} 
$M_y = 2x$, $N_x = 2x$. Exact.
$F(x,y) = \int M dx = x^2 y + x + g(y)$.
Differentiate wrt $y$ to find $g(y)$:
$F_y = x^2 + g(y)$. Set this equal to $N$:
$x^2 + g(y) = x^2 + \cos y \implies g(y) = \cos y$.
Solution: $x^2 y + x + \cos y = C$.
\newline \textbf{Find the Flaw:}
\vspace{1cm}

\noindent \textbf{Problem 4: The Expression Trap}
\newline \textbf{Equation:} $3x^2 dx + 3y^2 dy = 0$
\newline \textbf{Flawed Solution:} 
$M = 3x^2, N = 3y^2$. $M_y = 0$, $N_x = 0$. Exact.
$F(x,y) = \int 3x^2 dx = x^3 + g(y)$.
$F_y = g'(y) = 3y^2 \implies g(y) = y^3$.
The implicit solution to the differential equation is $x^3 + y^3$.
\newline \textbf{Find the Flaw:}
\vspace{1cm}

\noindent \textbf{Problem 5: The IVP Sign Trap}
\newline \textbf{Equation:} $2x dx + 2y dy = 0, \quad y(0) = -3$
\newline \textbf{Flawed Solution:} 
$M_y = 0, N_x = 0$. Exact.
$F(x,y) = x^2 + y^2 = C$.
Substitute $x=0, y=-3 \implies 0^2 + (-3)^2 = C \implies 9 = C$.
The equation is $x^2 + y^2 = 9$.
Solving for $y$, we get the explicit solution $y = \sqrt{9 - x^2}$.
\newline \textbf{Find the Flaw:}
\vspace{1cm}

\newpage
\subsection*{Answer Key for "Find the Flaw"}
\textbf{Flaw 1:} The student incorrectly added the overlapping term $x^2 y$ twice. When combining the integrals of $M$ and $N$, identical terms shared between them only appear \textit{once} in the potential function. The correct solution is $x^2 y + y^3 = C$.

\textbf{Flaw 2:} The student tested for exactness by taking the derivative of $M$ with respect to $x$, and $N$ with respect to $y$. The correct exactness test is the opposite: $M_y = N_x$. Because $M_y = 2$ and $N_x = 2$, the equation \textit{is} exact.

\textbf{Flaw 3:} When differentiating $F(x,y)$ with respect to $y$, the student wrote the derivative of $g(y)$ as $g(y)$ instead of $g'(y)$. Thus, they skipped the step of integrating $\cos y$. The correct step is $g'(y) = \cos y \implies g(y) = \sin y$. The correct solution is $x^2 y + x + \sin y = C$.

\textbf{Flaw 4:} The student left the final answer as an expression, not an equation. Exact differential equations result in an implicit equation set equal to an arbitrary constant. The correct final answer must be written as $x^3 + y^3 = C$.

\textbf{Flaw 5:} When taking the square root of $y^2 = 9 - x^2$, the student defaulted to the positive root. However, the initial condition requires $y(0) = -3$, which is negative. The positive root yields $y(0) = 3$. The correct explicit solution must use the negative branch: $y = -\sqrt{9 - x^2}$.

\end{document}